\documentclass[10pt,a4paper]{amsart}
\usepackage[margin=19mm]{geometry}
\usepackage{amsmath,amssymb,mathtools}
\usepackage[scr=rsfs,cal=euler]{mathalfa}
\usepackage{xfrac}
\usepackage[unicode,hidelinks,bookmarks=false]{hyperref}
\newcommand{\ie}{i.e.\ }
\newcommand{\cf}{cf.\ }
\newcommand{\resp}{respectively\ }
\newcommand{\Subsection}{\subsection}
\providecommand{\nobreakdash}{\nobreak\hbox{-}}
\setlength{\emergencystretch}{2em}
\multlinegap0pt
\usepackage{mdframed}
\usepackage{mdframed}
\usepackage{mdframed}
\usepackage{mdframed}
\usepackage{mathtools}
\usepackage{algorithm}\usepackage{algpseudocode}
\usepackage{tcolorbox}
\usepackage{bm}
\usepackage{amssymb}
\usepackage{xcolor}
\usepackage[shortlabels]{enumitem}
\usepackage{tikz}
\usepackage{mdframed}
\newtheorem{proposition}{Proposition}
\usepackage{cleveref}
\crefname{proposition}{Proposition}{Propositions}
\newcommand{\cR}{\mathcal{R}}
\newcommand{\dom}[1]{\mathrm{dom}(#1)}
\title{A New Kernel Regularity Condition for Distributed Mirror Descent: Broader Coverage and Simpler Analysis}
\author{Junwen Qiu \and Ziyang Zeng${}^*$ \and Leilei Mei${}^*$ \and Junyu Zhang${}^*$}
\date{Source: arXiv:2603.12838v2 (CC BY 4.0)}
\begin{document}
\maketitle

\noindent\emph{Source and licence.} A New Kernel Regularity Condition for Distributed Mirror Descent: Broader Coverage and Simpler Analysis. Version arXiv:2603.12838v2 (CC BY 4.0), distributed by arXiv under the Creative Commons Attribution 4.0 International licence, http://creativecommons.org/licenses/by/4.0/, as recorded in the arXiv licence metadata for this exact version and shown on the abstract page https://arxiv.org/abs/2603.12838v2. Full licence notice: \texttt{licenses/arxiv-2603.12838-CC-BY-4.0.txt}.\par\smallskip

\noindent\textit{Editorial scope.} Counted unit: the article's proposition 1 (source0.tex lines 518--525). The article's complete original proof of it is delivered with it (source0.tex lines 530--558, one \texttt{proof} environment of 29 lines, which the article places away from the statement). No mathematics is written, repaired or re-derived: the statement and the proof are the article's own lines, in the article's own notation, and no retained line is substituted or re-flowed.  Retained uncounted as the source-local context the proof needs: lines 467--470.  Nothing was omitted from any retained span, so the package carries no omission record.  No retained line contains a citation, so the wrapper carries no bibliography and no bibliography entry.  No retained line contains a figure environment, an image-inclusion command or drawing code, so no image is delivered and the package declares no asset dependency.  Editorial formatting only, in addition: an AMS \texttt{amsart} preamble that restates the article's own declarations and macros used by the retained lines, namely the environments \texttt{proposition} on separate counters, together with the standard packages those lines need.  The retained environments print in this document as Proposition 1 in order of appearance, whereas the article prints the same items with the numbering of its own full document.  The complete native source of the article as distributed in the arXiv e-print for this exact version is bundled unchanged as \texttt{sources/196295/source0.tex}.
\begin{equation}
    \label{eqn:HRUC-Hessian-comparison}
    \frac{1}{1+\zeta(\delta)}\nabla^2h(y)\preceq \nabla^2h(x)\preceq \big(1+\zeta(\delta)\big)\nabla^2h(y)\qquad \text{when}\qquad \rho_h(x,y)\leq\delta.
\end{equation} 

\begin{proposition}[\textbf{Closedness under Combination}]
    \label{proposition:Affine 1} 
    Let the kernels $h$ and $g$ be $\zeta_h$- and $\zeta_g$-HRUC regular, respectively. Let $\varphi:= h+ g$, then it holds that  \vspace{0.1cm}\\
    \noindent\emph{(i).} $\,\,\,\varphi$ is $\zeta_h$-HRUC if $g(x) = \|x\|^2/2$.\vspace{0.1cm}\\
    \noindent\emph{(ii).} $\,\varphi$ is $\max\{\zeta_h,\zeta_g\}$-HRUC if both $h$ and $g$ are coordinate-separable.\vspace{0.1cm}\\
    \noindent\emph{(iii).} $\!\varphi$ is $\max\{\zeta_h,\zeta_g\}$-HRUC if the kernel pair $(h,g)$ is cross-monotone in the sense that 
    $\langle \nabla h(x)-\nabla h(y),  \nabla g(x)-\nabla g(y)\rangle \geq 0$,  for $x,y\in\dom{\nabla\varphi}$.
\end{proposition}

\begin{proof}[Proof of \Cref{proposition:Affine 1}] First of all, through direct computation, we obtain 
\begin{equation*} 
\rho_\varphi^2(x,y) = \rho_h^2(x,y) + \rho_g^2(x,y) + 2\langle \nabla h(x)-\nabla h(y),  \nabla g(x)-\nabla g(y)\rangle. \end{equation*}
A direct consequence of the cross-monotonicity is that $\max\left\{\rho_h(x,y),\rho_g(x,y)\right\}\leq \delta$ when $\rho_\varphi(x,y)\leq\delta$. For example, in case (i) where $g$ is quadratic, and in case (ii) where both $h$ and $g$ are coordinate-separable, it is straightforward to verify the cross-monotonicity. Next, let us prove the three cases given $\max\left\{\rho_h(x,y),\rho_g(x,y)\right\}\leq \delta$.\vspace{0.1cm} 

\noindent\textbf{Case (ii).} $\,\,$ In the coordinate-separable case, denote $\nabla^2f(z) = \mathrm{Diag}(f_z^i)\succ0$ for $i\in[d]$, $z\in\{x,y\}$, and $f\in\{h,g\}$.  By cross-monotonicity,   
$\max_i\left\{\left|h_x^i/h_y^i-1\right|\right\}\leq \zeta_h(\delta)$ and  $\max_i\left\{\left|g_x^i/g_y^i-1\right|\right\}\leq \zeta_g(\delta)$. Then direct computation gives
\begin{equation*}
\begin{aligned}
     \cR_\varphi(x,y) &= \max_{i\in[d]}\left\{\left|\frac{h_x^i+g_x^i}{h_y^i+g_y^i}-1\right|\right\}  \\&= \max_{i\in[d]}\left\{\left|\frac{h_y^i(h_x^i/h_y^i-1)}{h_y^i+g_y^i}+\frac{g_y^i(g_x^i/g_y^i-1)}{h_y^i+g_y^i}\right|\right\} \leq \max\{\zeta_h(\delta),\zeta_g(\delta)\}.
\end{aligned}   
\end{equation*}
\noindent\textbf{Case (iii).} $\,\,$ Denote $H_z = \nabla^2h(z)\succ0$ and $G_z = \nabla^2g(z)\succ0$ for $z\in\{x,y\}$. Let $\zeta(\delta):=\max\{\zeta_h(\delta),\zeta_g(\delta)\}$. Applying the Hessian comparison \eqref{eqn:HRUC-Hessian-comparison}, we obtain
\begin{equation*}
\frac{1}{1+\zeta(\delta)}(H_y+G_y)
\preceq H_x+G_x
\preceq\big(1+\zeta(\delta)\big)(H_y+G_y).
\end{equation*}
Therefore, eigenvalues of $(H_y+G_y)^{-\frac12}(H_x+G_x)(H_y+G_y)^{-\frac12}$ lie in $[(1+\zeta(\delta))^{-1},1+\zeta(\delta)]$. Hence,
\begin{equation*}
\begin{aligned}
     \cR_\varphi(x,y)
     &=\left\|(H_y+G_y)^{-\frac12}(H_x+G_x)(H_y+G_y)^{-\frac12}-I\right\|\\
     &\leq \max\left\{\zeta(\delta),1-\frac{1}{1+\zeta(\delta)}\right\}
     =\zeta(\delta).
\end{aligned}
\end{equation*}
\noindent\textbf{Case (i).} $\,\,$ Finally, using the above bound with $\zeta_g(\delta)=0$ (as the quadratic kernel is HRUC with modulus $\zeta_g(\delta)=0$), we obtain $\cR_\varphi(x,y)\leq \zeta_h(\delta)$.
\end{proof}

\end{document}
